\origpage{33}{}
\textbf{Proof.} Let $H_j = \mu(F_j^*)$ ($1 \le j \le r$). We must check that (V1)--(V4) hold (and that $\mu(s) \subseteq H_j$ if $s \subseteq F_j^*$ is a space, which is obvious). For (V1),
\begin{equation*}
G = \Gamma \cup \bigcup(\mu(F_j^*) : 1 \le j \le r) = \Gamma \cup H_1 \cup \dots \cup H_r
\end{equation*}
from (P1). To see that $H_1, \dots, H_r$ are mutually disjoint, suppose that $v \in V(H_{j_1} \cap H_{j_2})$ say. For $i = 1, 2$ let $x_i$ be a unit in $F_{j_i}^*$ with $v \in V(\mu(x_i))$. From (P1), $x_1$ and $x_2$ belong to the same cuff $C$. We may assume that $k \le r$, and that for $1 \le j \le r, F_j^* \subseteq C$ if and only if $j \le k$, and that $F_1^*, \dots, F_k^*$ occur in order around $C$. Let $s_1^*, \dots, s_k^*$ be the spaces in $C$ not in $F_1^* \cup \dots \cup F_r^*$, numbered so that if $k > 1$ then $s_j^*$ has one end in $F_j^*$ and one in $F_{j+1}^*$ for $1 \le j < k$, and $s_k^*$ has one end in $F_k^*$ and one in $F_1^*$. For $1 \le j \le k$, let $F_j^*$ have ends $a_j, b_j$, where $a_j$ is an end of $s_j^*$ and $b_j$ is an end of $s_{j-1}^*$ (or of $s_k^*$, if $j = 1$). Let us write $b_{k+1} = b_1$ and $a_0 = a_k$. Now since the pseudo-presentation is disjointed, $v$ does not belong to both $\mu(a_{j_1-1})$ and $\mu(b_{j_1})$; choose $u_1 \in \{a_{j_1-1}, b_{j_1}\}$ so that $v \notin V(\mu(u_1))$. Similarly choose $u_2 \in \{a_{j_1}, b_{j_1+1}\}$ with $v \notin V(\mu(u_2))$. Now $x_1, u_1, x_2, u_2$ occur in $C$ in order, contrary to (P3), since $v \in V(\mu(x_1) \cap \mu(x_2))$ and $v \notin V(\mu(u_1) \cup \mu(u_2))$.

This proves that $H_1, \dots, H_r$ are disjoint. To complete the proof of (V1), we must check that for $1 \le j \le r, V(H_j \cap \Gamma) = V(\Gamma) \cap F_j^*$, and $E(H_j \cap \Gamma) = \emptyset$. Now $V(\Gamma) \cap F_j^* \subseteq V(\mu(F_j^*)) = V(H_j)$, because $u \in V(\mu(u))$ for each $u \in V(\Gamma) \cap F_j^*$. Conversely, let $u \in V(H_j \cap \Gamma)$. Let $x$ be a unit in $F_j^*$ with $u \in V(\mu(x))$. Since $V(\mu(x) \cap \Gamma) = \{x\}$ if $x$ is a vertex, and $V(\mu(x) \cap \Gamma)$ is the set of ends of $x$ by (P2) if $x$ is a space, it follows that either $\delm{v}\fix{u} \in V(\Gamma) \cap F_j^*$ or $\delm{v}\fix{u}$ is an end of some space in $F_j^*$; and so in either case, $\delm{v}\fix{u} \in V(\Gamma) \cap F_j^*$. This proves that $V(H_j \cap \Gamma) = V(\Gamma) \cap F_j^*$. Finally, $E(H_j \cap \Gamma) = \emptyset$ because $E(\Gamma \cap \mu(\delm{S}\fix{s})) = \emptyset$ for each space $s \subseteq F_j^*$, by (P1). This proves (V1).\fixnote{fix:8.1-proof-V1}{The vertex in this sentence is $u$, chosen just before, and $v$ names the vertex of the paragraph before. Also $\mu(S)$ is $\mu(s)$, for the space $s$.}

Now (V2) holds by definition, and (V3) is the same as (P2), and so it only remains to verify (V4). Let $1 \le j \le r$, and let $s_1, s_2 \subseteq F_j^*$ be spaces, and let $u \in F_j^* \cap V(\Gamma)$ lie in $F_j^*$ between $s_1$ and $s_2$. Let $s_0$ be a space with $s_0 \nsubseteq F_1^* \cup \dots \cup F_r^*$, in the same cuff $C$ as $F_j^*$. (This exists since $|V(\Gamma) \cap C| \ge 2$.) Let $s_0$ have ends $a_0$ and $b_0$. Now $a_0, s_1, u, s_2$ occur in order in $C$, and so by (P3), $\mu(s_1) \cap \mu(s_2) \subseteq \mu(a_0) \cup \delm{\mu(s_2)}\fix{\mu(u)}$. Similarly, $\mu(s_1) \cap \mu(s_2) \subseteq \mu(b_0) \cup \delm{\mu(s_2)}\fix{\mu(u)}$, and so
\begin{equation*}
\mu(s_1) \cap \mu(s_2) \subseteq (\mu(a_0) \cap \mu(b_0)) \delm{\cap}\fix{\cup} \delm{\mu(s_2)}\fix{\mu(u)} = \delm{\mu(s_2)}\fix{\mu(u)}.\fixmark{fix:8.1-proof}
\end{equation*}
\fixnotetext{A misprint: (V4) asks for $\mu(s_1) \cap \mu(s_2) \subseteq \mu(u)$, and (P3), applied to $a_0, s_1, u, s_2$, gives $\mu(s_1) \cap \mu(s_2) \subseteq \mu(a_0) \cup \mu(u)$. So $\mu(s_2)$ on the right is $\mu(u)$, all four times, and the second $\cap$ in the display is $\cup$.}%
\fix{Here $\mu(a_0) \cap \mu(b_0)$ is null, since the pseudo-presentation is disjointed and $s_0 \nsubseteq F_1^* \cup \dots \cup F_r^*$.} This proves (V4), as required. $\square$

Secondly, we wish to discuss splitting vertices in a pseudo-presentation, in order to make it disjointed. More precisely, let $\Gamma, F_1^*, \dots, F_r^*, \mu$ be a pseudo-presentation of $G$ in $\Sigma$. Let $s^*$ be a space with $s^* \nsubseteq F_1^* \cup \dots \cup F_r^*$, with ends $u_1, u_2$ say. Let $C$ be the cuff including $s^*$. Let $v \in V(\mu(u_1) \cap \mu(u_2))$.
